\section{Experiment 2: The Boundaries of Perturbation}\label{sec:exp2}

\textbf{Objective:} Map the information-geometric walls that constrain where equilibria can exist. Homeostasis only operates where the system has degrees of freedom.

\subsection{Observable-Dependent Saturation}

\paragraph{Hypothesis.}
Perturbing the equilibrium through architectural or training interventions should improve performance in sub-saturated regimes (where the system has representational capacity to spare). In saturated regimes (where information channels are at capacity), the system should be rigid and unresponsive to perturbation.

\paragraph{Methods.}
We systematically vary perturbation magnitude $\omega \in [0.5, 1.5]$ on suppressor head weights and measure two observables:
\begin{enumerate}
\item \textbf{Power:} Grokking speed (steps to 95\% test accuracy)
\item \textbf{Information:} Mutual information (MI) between layer activations, estimated via k-NN density estimation~\cite{kraskov2004estimating}
\end{enumerate}

\textbf{Operational Definitions.} MI is computed between consecutive layer representations: $I(h^{(l)}; h^{(l+1)})$ where $h^{(l)}$ are activation vectors at layer $l$. Effective dimensionality $d_{\text{eff}}$ is measured via intrinsic dimensionality estimation on embedding spaces. Both metrics are task- and architecture-dependent; the thresholds reported (MI $\approx$ 2--3 bits, $d_{\text{eff}} \approx$ 30--50) reflect empirical saturation points observed across our experiments and are consistent with prior work on representational capacity limits in transformers.

We test this across tasks of varying difficulty to identify where information saturation occurs.

\paragraph{Results: The U-Shaped Omega Curve.}
In sub-saturated regimes (simple tasks with high representational slack), perturbations in either direction ($\omega < 1$ or $\omega > 1$) accelerate grokking. The system exhibits a U-shaped curve with minimum at $\omega = 1.0$ (the natural equilibrium). Destabilizing the equilibrium provides "activation energy" to escape memorization plateaus.

However, as task difficulty increases and information channels approach saturation, this effect disappears. Performance plateaus regardless of $\omega$, and the U-shaped curve flattens.

\paragraph{The Saturation Boundary.}
We identify a critical threshold where observable-dependent metrics saturate:
\begin{itemize}
\item \textbf{Below threshold:} System is plastic, responds to perturbation, U-shaped curve persists
\item \textbf{At threshold:} Mutual information $\approx$ 2--3 bits, channel capacity exhausted
\item \textbf{Above threshold:} System is rigid, ignores perturbation, flat response curve
\end{itemize}

\subsection{Capacity-Contingent Attractor Dynamics}

The equilibrium's functional role diverges qualitatively with model scale, revealing capacity-dependent phase transitions.

\paragraph{Small Scale (GPT-2 Small, 124M parameters, $d_{\text{eff}} \approx 15$).}
Individual ablations of suppressor heads on high-uncertainty tasks \emph{reduce} logit difference by 1.1--5.5\% relative to baseline:
\begin{itemize}
\item Negation: head 0:2 $\Delta$LD = $-0.042$, head 0:4 = $-0.101$, head 0:7 = $-0.032$
\item Counterfactual: head 0:2 = $-0.030$, head 0:4 = $-0.039$, head 0:7 = $-0.017$
\item Logic: head 0:2 = $-0.054$, head 0:4 = $-0.067$, head 0:7 = $-0.019$
\end{itemize}

The triplet ablation shows identical negative direction (mean $\Delta$LD: negation $-0.150$ [$-0.202, -0.111$], counterfactual $-0.081$, logic $-0.097$ [$-0.130, -0.057$]). Output-entropy deltas are small and positive ($\Delta H = +0.03$ to $+0.34$ nats).

\textbf{Interpretation:} At small scale, suppressors act as \emph{capacity-consuming gates} rather than global attractor-setters. Ablation can improve factuality by forcing the model to reallocate limited capacity toward the task.

\paragraph{Medium Scale (GPT-2 Medium, 355M parameters, $d_{\text{eff}} \approx 50$--80).}
Identical ablations \emph{improve} logit difference by 40--85\% across all probes (mean $\Delta$LD: facts $+0.40$, negation $+0.84$, counterfactual $+0.85$, logic $+0.55$). H6 path mediation is robust ($\approx 67\%$ on facts), and the geometry signature is strong (output-entropy $\Delta H = -2.4$ to $-3.8$ nats).

\textbf{Interpretation:} Above the capacity threshold, suppressors stabilize a global hedging attractor that downstream layers inherit and reinforce. The bias toward uncertainty becomes a stable system-wide feature.

\subsection{The "Walls" of the Manifold}

The equilibrium manifold is bounded by information-theoretic constraints:
\begin{enumerate}
\item \textbf{Sub-saturated region (plastic):} System responds to perturbation, can explore different equilibria, U-shaped omega curve visible
\item \textbf{Saturation boundary (the wall):} MI $\approx$ 2--3 bits, channel capacity exhausted, $d_{\text{eff}} \approx 30$--50
\item \textbf{Saturated region (rigid):} System ignores perturbation, equilibrium frozen, flat response curve
\end{enumerate}

\begin{table}[h]
\centering
\caption{Cross-scale summary: capacity-contingent attractor dynamics.}
\label{tab:capacity-boundaries}
\begin{tabular}{lccc}
\toprule
Metric & Small (124M) & Medium (355M) & Interpretation \\
\midrule
L0 triplet $\Delta$LD (neg) & $-0.150$ & $+0.842$ & Sign flip: maladaptive $\to$ adaptive \\
L0 triplet $\Delta$LD (logic) & $-0.097$ & $+0.552$ & Sign flip: maladaptive $\to$ adaptive \\
H6 mediation (facts) & 5.9\% & 67\% $\pm$ 4\% & Attractor basin only at scale \\
Output entropy $\Delta H$ & $+0.03$ to $+0.34$ & $-2.4$ to $-3.8$ & Geometry scales with capacity \\
Interpretation & Local gates & Global attractor & Phase-transition regime \\
\bottomrule
\end{tabular}
\end{table}

\subsection{Findings Summary}

\begin{enumerate}
\item \textbf{Boundary exists:} Information saturation creates a hard wall at MI $\approx$ 2--3 bits
\item \textbf{Below boundary:} System is plastic, perturbations accelerate learning
\item \textbf{At boundary:} System becomes rigid, perturbations have no effect
\item \textbf{Scale dependence:} Suppressor role transitions from maladaptive (small scale) to adaptive (large scale) at capacity threshold
\item \textbf{Implication:} Equilibria can only exist where the system has representational degrees of freedom
\end{enumerate}

\textbf{Interpretation:} This establishes the boundaries for Experiments 3 and 4. Homeostatic compensation can only occur in sub-saturated regimes. When the system hits the "wall" of channel capacity, the equilibrium becomes frozen. The space of possible equilibria is fundamentally constrained by the information geometry of the architecture.
